
Hallucination detection methods rely on Natural Language Inference (NLI) models to assess claim-context consistency. However, NLI models trained on SNLI/MNLI (semantic similarity tasks) may not generalize to factual verification. This raises competing explanations for detection failures: does the method fail because creative text confuses the model (domain-specific hypothesis), or because the NLI component was never properly calibrated for factual verification (measurement validity issue)?

Claim-Conditioned Probability (CCP) aggregates token-level probabilities weighted by NLI-derived entailment status, reporting +0.05--0.10 ROC-AUC improvements over baselines. We set out to test whether CCP degrades when applied to creative text---fiction, poetry, metaphorical content---compared to factual text. Our hypothesis: CCP's NLI-based conditioning embeds implicit factual-ontology assumptions (e.g., claims must correspond to verifiable facts) that misalign with creative semantics, where metaphors and speculation are legitimate rather than erroneous. We predicted that the claim-type mass ratio ($\rho_j$, the core CCP diagnostic metric) would drop by >0.15 when applied to creative vs factual text.

Instead, we could not reproduce the baseline. Across both factual text (TruthfulQA) and creative text (WritingPrompts), $\rho_j$ values were 20-$80\times$ lower than the inferred range: median 0.0354 (factual) and 0.0103 (creative) vs the inferred range of 0.75--0.85 from the CCP paper's ROC-AUC claims. Statistical tests showed no domain separation ($p = 1.0$, Cohen's $d = -0.0635$).

Root cause analysis revealed that DeBERTa-v3-base, the NLI model trained on SNLI/MNLI (sentence-pair semantic similarity tasks), assigns approximately 90\% probability mass to the ''neutral'' class for claim-context pairs in both domains. This is not a domain-specific failure (creative text confusing the model) but a case study of known NLI miscalibration for factual verification: SNLI/MNLI training objectives optimize for semantic similarity detection (''Do these sentences describe similar situations?''), not factual verification (''Is this claim consistent with this context?''). The model defaults to ''neutral'' for claim-context pairs with limited lexical overlap, mechanistically driving $\rho_j$ toward zero.

This failure mode teaches a critical methodological lesson: measurement validity is prerequisite for hypothesis testing. When a metric produces values 20-$80\times$ outside the inferred range across all conditions, you cannot distinguish ''hypothesis is wrong'' from ''measurement is broken.'' We faced a logical impossibility: is the uniform degradation evidence that creative text does not confuse CCP (hypothesis refuted), or that our CCP implementation does not work as described (measurement broken)? Without baseline replication on the original domain, we cannot separate these explanations.

\textbf{Contributions}:

\begin{enumerate}
\item \textbf{Transparent Failure Documentation}: We provide the first systematic replication attempt of CCP, documenting both what went wrong and why. This negative result is itself a contribution---it exposes a reproducibility gap in the hallucination detection literature.
\end{enumerate}

\begin{enumerate}
\item \textbf{Root Cause Hierarchy}: We identify NLI calibration (Tier 1: primary), claim decomposition quality (Tier 2: contributory), and context pairing strategy (Tier 2: contributory) as the failure modes, with evidence strength rankings and falsifiability tests for each.
\end{enumerate}

\begin{enumerate}
\item \textbf{Methodological Requirements}: We adapt Dodge et al. (2019) reproducibility checklist for hallucination detection: (R1) report raw metric distributions, not just ROC-AUC; (R2) validate NLI calibration on known examples before running experiments; (R3) document claim decomposition methodology with inter-method agreement; (R4) provide public code with baseline replication notebooks.
\end{enumerate}

\begin{enumerate}
\item \textbf{Case Study of NLI Miscalibration}: We provide a case study illustrating that NLI miscalibration for factual verification---a known issue in domain adaptation literature---can prevent hypothesis testing in hallucination detection research.
\end{enumerate}

\textbf{Broader Impact}: With 50+ hallucination detection papers published in 2024 citing NLI-based methods, transparent failures accelerate progress by preventing repetition of costly mistakes. Our findings suggest that hallucination detection papers must adopt higher standards for implementation transparency---raw metric distributions, NLI calibration diagnostics, and claim decomposition validation should be first-class contributions, not afterthoughts.

The remainder of the paper is organized as follows. Section 2 reviews hallucination detection methods, NLI domain adaptation, and reproducibility challenges. Section 3 documents our CCP implementation with full transparency. Section 4 presents experimental results, including the gate failure and NLI distribution analysis. Section 5 performs root cause analysis via competing explanations framework. Section 6 discusses broader implications, limitations, and future work. Section 7 concludes with lessons learned and recommendations for the field.
